% command to make sure colorboxes don't create whitespace between code lines
% Reduce the space inside colorboxes to avoid extra whitespace between lines
\fboxsep=.6pt

% default settings for listings
\lstset{style=python}

\title[Daten Vorhersagen]{Daten Vorhersagen}

\subtitle{Markov-Ketten}

\author[Wendl, Cyril] % (optional)
{Naoki Peter\inst{1} \and Cyril Wendl \inst{2}}

\institute[KZN] % (optional)
{
    \inst{1}%
    Fachschaft Informatik\\
    Kantonsschule Zürich-Nord
    \and
    \inst{2}%
    Fachschaft Informatik\\
    Kantonsschule im Lee
}
\date[\today] % (optional)

\logo{\includegraphics[width=\logosize]{\thelogo}}

\titlegraphic{\vspace{-.7cm}\includegraphics[width=\textwidth]{"GrundlagenInfo/06_Datenbanken/Figures/title-emojis"}}

\frame{\titlepage}


\begin{frame}{Themenübersicht}
    \begin{enumerate}
        \item[\faCheckSquare{}] Daten verstehen (analysieren und visualisieren)
        \item[\faCheckSquare{}] Zahlen vorhersagen
        \item[\faSquare] \textbf{Kategorien vorhersagen}
    \end{enumerate}
\end{frame}

\begin{frame}{Maschinelles Lernen}
    \begin{figure}[H]
        \only<1->{\centering\includegraphics[width=0.26\pagewidth]{"GrundlagenInfo/10_AusDatenLernen/Figures/intro-markov-weather"}}
        \only<2->{\centering\includegraphics[width=0.26\pagewidth]{"GrundlagenInfo/10_AusDatenLernen/Figures/intro_chatGPT-internals.jpeg"}}
        % \only<3->{\centering\includegraphics[width=0.26\pagewidth]{"GrundlagenInfo/10_AusDatenLernen/Figures/intro_chatGPT-internals"}}
    \end{figure}
\end{frame}

\begin{frame}{Kategorien vorhersagen}{Beispiel: Wetter}
    Statistisch beobachtete Wetterdaten zeigen, dass das Wetter in der Regel nicht zufällig ist:
    \begin{itemize}
        \item Wenn es heute regnet, ist die Wahrscheinlichkeit hoch, dass es morgen auch regnet.
        \item Wenn es heute sonnig ist, ist die Wahrscheinlichkeit hoch, dass es morgen auch sonnig ist.
    \end{itemize}
    Diese Daten können gemessen und als Übergangstabelle dargestellt werden:

    \begin{table}[h]
        \centering
        \begin{tabular}{l|c|c}
            & morgen sonnig & morgen regnerisch \\
            \hline
            heute sonnig     & 0,8 & 0,2 \\
            heute regnerisch & 0,4 & 0,6 \\
        \end{tabular}
    \end{table}
\end{frame}


\begin{frame}{Auftrag}{Jupyter Notebooks}
    \begin{itemize}
        \item Aufgabenblatt 3a durchlesen und vervollständigen
        \item Aufgabenblatt 3b lösen
    \end{itemize}
\end{frame}